\documentclass[11pt,a4paper]{article}
\usepackage[utf8]{inputenc}
\usepackage[spanish]{babel}
\usepackage{amsmath,amsfonts,amssymb}
\usepackage{geometry}
\geometry{top=2.5cm, bottom=2.5cm, left=2.5cm, right=2.5cm}

\begin{document}

\begin{center}
    \textbf{\Large UNIVERSIDAD CATÓLICA BOLIVIANA ``SAN PABLO'' -- SEDE TARIJA}\\
    \vspace{0.2cm}
    \textbf{Carrera: Ingeniería Mecatrónica | Asignatura: Robótica (IMT-342)} \\ %[cite: 1]
    \textbf{Docente:} M.Sc. Ing. Bernardo Quiroga Turdera
\end{center}

\vspace{0.3cm}
\noindent
\textbf{Estudiante:} \rule{8cm}{0.4pt} \hfill \textbf{Fecha:} 04/08/2026 \\

\noindent
\textit{Instrucciones: Desarrolle de forma manuscrita los siguientes ejercicios. Esta prueba es formativa (ponderación 0\%) y busca identificar fortalezas y áreas de reforzamiento en matemática y control. Tiempo: 45 minutos.}

\vspace{0.5cm}
\noindent \textbf{Parte A: Álgebra Lineal y Geometría Espacial (40\%)}

\begin{enumerate}
    \item Dado un vector en el plano $p = [2, 3]^T$, obtenga la matriz de rotación $R \in \mathbb{R}^{2\times2}$ para un giro de $\theta = 30^\circ$ en sentido antihorario respecto al origen y calcule el vector resultante $p'$.
    
    \vspace{2.5cm}
    \item Si $R \in \mathbb{R}^{3\times3}$ es una matriz de rotación pura en el espacio, demuestre por definición por qué se cumple estrictamente que $\det(R) = +1$ y que su inversa equivale a su transpuesta ($R^{-1} = R^T$).
    
    \vspace{2.5cm}
    \item Explique brevemente qué representa físicamente la Descomposición en Valores Singulares (SVD) de una matriz de datos reales.
    \vspace{2cm}
\end{enumerate}

\noindent \textbf{Parte B: Teoría de Control Clásico (40\%)}

\begin{enumerate}
    \setcounter{enumi}{3}
    \item Un sistema físico dinámico está modelado por la ecuación diferencial:
    $$ \ddot{\theta}(t) + 5\dot{\theta}(t) + 4\theta(t) = \tau(t) $$
    Obtenga la función de transferencia $G(s) = \Theta(s)/\tau(s)$ en el dominio de Laplace y determine si el sistema es sobreamortiguado, subamortiguado o críticamente amortiguado.
    
    \vspace{3cm}
    \item Dibuje el diagrama de bloques de un lazo de control PID en lazo cerrado con realimentación negativa y explique qué efecto físico tiene un exceso de ganancia derivativa ($K_d$) sobre el ruido de un sensor.
    \vspace{3cm}
\end{enumerate}

\noindent \textbf{Parte C: Lógica y Programación (20\%)}

\begin{enumerate}
    \setcounter{enumi}{5}
    \item Describa en pseudocódigo o código de Python un bucle anidado (\texttt{for}) para transponer una matriz de $3 \times 3$ sin usar funciones directas de NumPy.
    
    \vspace{3cm}
    \item ¿Cuál es la diferencia entre un proceso de ejecución síncrono (blocking) y uno asíncrono gestionado por eventos en sistemas de tiempo real?
\end{enumerate}

\end{document}